\documentclass[14pt,a4paper]{article}
\usepackage[utf8]{inputenc}
\usepackage[paper=a4paper, margin=1.5cm]{geometry}

% FA icons
\usepackage{hyperref}
\usepackage{url}
\usepackage{color}
\usepackage{graphicx}

\usepackage{amsmath}
\usepackage{amssymb}

\hypersetup{
  colorlinks=false,% hyperlinks will be black
  pdfborderstyle={/S/U/W .5}% border style will be underline of width 1pt
}

% citace a jazyk
\usepackage[main=czech, english]{babel}
\usepackage{csquotes}
\usepackage[style=iso-authoryear]{biblatex}
\addbibresource{zdroje.bib}

\begin{document}
	\begin{titlepage}
		\begin{center}
            {
            \centering
            \includegraphics[]{./PDF_tisk/UP_logo_PdF-UP_horizont_cz.pdf}
            }
			
			\vspace{3cm}

            {
                \LARGE
                \textbf{Šimon Janča}\\
                3.\,ročník\\[8mm]
                Obor: Informační technologie a Matematika pro vzdělávání
            }

            \vspace{4cm}
			
			{
			    \textbf{\Huge Využití IT ke generování zadání písemných testů}\\[4mm]
			    \Large
			    Bakalářská práce
			}

            \vfill
            
            {
                Vedoucí práce:
                doc. RNDr. Petr Šaloun, Ph.D.
                \hfill
    			Olomouc \the\year{}
            }
			
		\end{center}
	\end{titlepage}
	% generate table of contents
	\tableofcontents
	
	\newpage
	
	\section{Úvod}
    V~této práci si~kladu za~cíl vytvořit webovou aplikaci, kterou by~učitelé mohli používat k~tvorbě různých variant zadání písemných prací.
    
    Zároveň bych rád popsal základní prvky vývoje moderní \emph{MVC} aplikace a~popsal důvody k~jednotlivým rozhodnutím od~návrhu
    a~výběru technologií, přes plánování pomocí diagramů, až k~její implementaci a spuštění za~použití vybraných technologií.
    
    Pochopitelně nepůjde o~jediný způsob, jak mohou aplikace vzniknout, nicméně by~práce měla být čitelným a~reprodukovatelným postupem vývoje
    a~údržby podobné aplikace.
        
        \subsection{V současnosti používané nástroje na generování prací}
        % TODO: 
        \subsection{Co používají učitelé}

        \subsection{Účel a~popis fungování aplikace}
        Aplikace \emph{Písemkomat}, poskytne registrovaným uživatelům možnost vytvořit vytvořit si příklady, které může umístit do~volitelně
        strukturovaných kategorií - podle předmětu, ročníku, tématu... Tyto příklady může následně využít při tvorbě písemných prací,
        kde si bude moci zvolit, kategorie ze kterých se příklady vyberou a~počet příkladů, které se~mají do~písemky vygenerovat. 
        \subsection{Příklad použití}
        Mějme prvního uživatele, říkejme mu~třeba Adam. Adam je~učitel a~chce vytvořit 2~verze písemky s~rozdílnými příklady pro~své žáky.
        Adam se~zaregistruje do~aplikace \emph{Písemkomat} a~přihlásí se~do~svého účtu. Adam učí Matematiku, takže si~vytvoří 5 příkladů z~tohoto předmětu
        a~vloží je do kategorie Matematika.

        Následně si vytvoří písemnou práci a díky tomu, že si~Adam vytvořil 5~příkladů, může si~vybrat, zda chce,
        aby se~v~jeho písemce objevily všechny příklady, nebo jen některé. Bude tedy požadovat 3 příklady s~náhodně dosazenými hodnotami.
        Některé hodnoty nemohou být zcela náhodné, protože by~to mohlo vést k~neřešitelným příkladům, nebo příkladům s~nulovým výsledkem.
        Proto \emph{Písemkomat} umožní vytvořit si~proměnné, které budou mít přiřazenou hodnotu z~určitého intervalu nebo vypočtené z~jiných proměnných.

        Adam si~např.~vytvoří proměnné $A, B, C$ a $D$, kde $A, B$ budou náhodná čísla od 1 do 10, $C = 4/3*A$ a $D = C + B$.

        Řešený příklad se~vygeneruje do vzorce, který posléze zapíše do~\emph{Písemkomatu}. Do proměnných $A, B$ se uloží náhodné číslo
        z~intervalu, který~si~Adam nastavil a~pomocí zápisu proměnné - \texttt{\$A} nebo formátu umožňující
        delší názvy\,--\,\texttt{\$\{nazevpromenne\}}\,--\,se~vloží do~příkladu na~zadané místo, resp.~do zápisu výpočtu dalších proměnných,
        kde budou použity ve výpočtu.
        
        Tedy v~tomto případě předpokládáme, že~budou vygenerována náhodná čísla $A = 3$ a~$B = 5$
        a~hodnoty proměnných $C, D$ se~dopočítají takto:
        \begin{align*}
            A &= 3 \\
            B &= 5 \\
            C &= 4/3*A = 4/3*3 = 4 \\
            D &= C + B = 4 + 5 = 9
        \end{align*}

        Proměnné se při generování výsledné podoby příkladu nahradí vygenerovanými čísly nebo výrazy.
        Zároveň bude možné je~nastavit tak, aby hodnoty nevyčíslovaly. Proměnné budou také moci nést jiné objekty, např.~pole pro~textový vstup
        nebo výběr z~možností, čehož bude možné využít při~tvorbě zadání, kde~bude~více vstupních polí, která bude mít za úkol žák vyplnit.

        Výsledná podoba příkladu jednoduchého výrazu může vypadat např.~takto:
        \begin{align*}
            \$Ax + \$B &= \$C + \$D \\
            3x + 5 &= 4 + 9
        \end{align*}

        Adam bude mít taky možnost zvolit~si, jestli chce otevřenou nebo uzavřenou formu příkladu s možnostmi výběru.
        V~tomto případě může taky využít proměnné, které si vytvořil dříve. V obou případech může taky zapsat vzorec,
        ze~kterého se~vypočtou možnosti a označit správnou možnost.
        Ve výsledném příkladu se~v~možnostech výběru proměnné nahradí a~výraz se~vyčíslí.

        Výsledná podoba zadání a~možností příkladu může vypadat např.~takto:
        \begin{itemize}
            \item $3x + 5 = 4 + 9$
            \item $x = 2$
            \item $x = 3$
            \item $x = 4$
            \item $x = 5$
        \end{itemize}

        Vstupní políčka se~ve~výsledné práci vykreslí jako prázdné, podtržené místo o~velikosti určené při~zadávání proměnné.
        Nebo v~případě možnosti výběru se~na~místě vykreslí všechny možnosti, aby~žák mohl vybrat správnou možnost kroužkováním nebo tak, jak učitel určí.
        
        To umožní systém používat např.~pro jazykové předměty, kde se~využívá doplňování slov do~textu nebo~výběr písmene.

        Všechna nastavení bude možné ovlivnit (vypínat, zapínat, měnit velikosti) i~na úrovni celých písemek.

    \section{Vývoj moderních webových aplikací}
    V dnešní době je~webová aplikace první volbou pro~většinu vývojářů. Je~to~díky tomu, že~webová aplikace je~dostupná z~jakéhokoliv zařízení,
    které má~přístup k~internetu a~je~schopné zobrazit webový prohlížeč. Webová aplikace je~také snadno aktualizovatelná a~není potřeba ji
    instalovat na~každém zařízení zvlášť. Vývojáři se~také nemusí starat o~kompatibilitu s~různými operačními systémy, protože webový prohlížeč
    je~dostupný na~většině z~nich. Webová aplikace je~také snadno škálovatelná, protože se~spouští na~serveru a~uživatelé se~k~ní připojují.
    Navíc je možné díky různým nástrojům aplikace snadno vydat z~jednoho zdrojového kódu jak webovou tak nativní (Android, IOS) aplikaci.

    Pokud se~chceme pustit do~vývoje webové aplikace, která bude využívat výhod moderních postupů a~technologií, je~potřeba
    si~ujasnit pár současných i~historických pojmů a~pochopit, jak~spolu souvisí.

        \subsection{Základní statická HTML stránka}
        V roce 1989 vytvořil Tim Berners-Lee první webový prohlížeč a~jazyk HTML. HTML je~zkratka pro~HyperText Markup Language (Hypertextový značkovací jazyk).
        Základem je HTML stránka~--~soubor, který obsahuje textový obsah a~HTML značky, které určují, význam jednotlivých částí obsahu.
        Navíc lze jednotlivé soubory propojit pomocí hypertextových odkazů, které umožňují uživateli kliknutím přecházet z~jedné stránky na~druhou.

        V dnešní době se~HTML stránky skládají z~HTML značek, ale i~CSS stylů, které určují jak má prohlížeč jednotlivé části stránky zobrazit.
        K tomu se přidává i~JavaScript, který umožňuje přidat do~stránky dynamický obsah a~interaktivitu. \parencite{berners:1989:proposal}

        Jednotlivé prvky je možné také označit pomocí identifikátorů a~tříd, které se~používají pro~jejich identifikaci.
        Zatímco tříd může mít prvek více, identifikátor může mít jen jeden. Identifikátor by měl být unikátní. V jedné HTML stránce by se měl vyskytovat jen jednou.
        Třídy se~používají pro~skupiny prvků, které mají stejný vzhled. Prvek může mít více tříd a~kombinovat tak jejich vlastnosti.

        \subsection{Kompilátor a Interpret}
        Pro~lepší pochopení programování aplikací je~potřeba si~ujasnit rozdíl mezi kompilovaným a~interpretovaným jazykem.

        \textbf{Kompilovaný} jazyk je~jazyk, který se~pomocí tzv.\,překladače (kompilátoru) překládá do~binárního souboru, který obsahuje přímo instrukce,
        které dokáže počítač (procesor) vykonávat. Při~každém spuštění se~do paměti načte binární soubor a~procesor vykonává jeho instrukce.

        \textbf{Interpretovaný} jazyk naopak vykonáván programem, kterému říkáme interpret. Ten čte příkazy a~přímo je~vykonává. To~znamená,
        že~při~každém spuštění se~příkazy překládají znovu. Interpret je program, který musí běžet, zatímco příkazy interpretuje,
        což je taky jeden z~důvodů, proč je~interpretovaný jazyk pomalejší a náročnější na využítí zdrojů, než~kompilovaný.
        
        Interpretování může ale přinést i~výhody. Například je~možné příkazy vykonávat ihned, bez~nutnosti překladu.
        
        \subsection{Javascript}
        Javascript je skriptovací jazyk, který se~používá k~tvorbě interaktivních webových stránek. Dokáže reagovat na~uživatelské akce,
        jako je~kliknutí, a~tedy na~rozdíl od~statické podoby HTML stránky dokáže měnit obsah podle toho, jak~uživatel se~stránkou pracuje.

        Vznikl v~roce 1995 a~jeho tvůrci jsou Brendan Eich a~společnost Netscape.
        Javascript je~dnes nejrozšířenějším programovacím jazykem na~světě a~používá se~nejen pro~webové stránky,
        ale i~pro~tvorbu nativních aplikací.

        Díky spouštěčům Javascriptu, kterým říkáme Javascriptové enginy, je~možné spouštět Javascript na~různých platformách.
        A tak můžeme Javascript spouštět v~prohlížeči, ale i~na~serveru, a dokonce na~mikrokontrolerech. Jeden jazyk tak zvládne pokrýt
        většinu potřeb, které můžeme při~vývoji aplikací mít.

        Javascript je interpretovaný jazyk, což znamená, že~jednotlivé příkazy se~vykonávají postupně a~při~každém spuštění programu se~překládají.
        To přináší výhodu v~tom, že je možné příkazy vykonávat ihned, ale zároveň to znamená, že~při~každém spuštění se~musí překládat,
        což může být pomalejší než spouštění kompilovaného binárního souboru.

        Javascript dokáže příkazy zpracovávat synchronně i~asynchronně.
        
        Synchronní zpracování znamená, že~se~příkazy vykonávají postupně, tak jak
        jsou zapsané.
        
        Asynchronní zpracování znamená, že~se~příkazy vykonávají nezávisle na~sobě. To je~užitečné např.~při~načítání dat ze~serveru,
        které může trvat delší dobu. V~tomto případě se~synchronní příkazy vykonávají dál a~když se~data úspěšně načtou, vykoná~se příkaz, který na~ně čeká.

        \subsection{Typescript}
        Další nevýhodou Javascriptu je, že~není typově bezpečný. To znamená, že~není možné při~překladu zkontrolovat, zda~jsou všechny proměnné
        a~funkce použity správně. To může vést k~chybám, které se~projeví až při~spuštění programu. To je~pro~větší projekty nežádoucí, proto se~používají
        jiné jazyky, které jsou typově bezpečné a~po~kontrole jsou přeloženy\footnote{kompilují, transpilují} zpět do Javascriptu. Jedním z~nich je~Typescript.
        Kromě typové bezpečnosti přináší Typescript i~další výhody a~optimalizace, které zjednodušují vývoj a~zvyšují výkon aplikace.

        Typescript přidává do~Javascriptu statickou kontrolu typů, což znamená, že~proměnné mají přiřazený datový typ, což
        u~Javascriptu normálně není nutné a~je~možné typy měnit za~běhu programu.
        Díky kontrole typů je~možné provést při překladu kontrolu, zda~jsou všechny proměnné a~funkce použity správně, což redukuje množství chyb.
        Typescript se~kompiluje (transpiluje) zpět do Javascriptu a~je~možné při~vývoji v~Typescriptu použít všechny vlastnosti Javascriptu.

        \subsection{DOM - Document Object Model}
        DOM je~objektový model dokumentu, který vytváří prohlížeč na~základě HTML stránky. Jednotlivé HTML značky jsou reprezentovány objekty.
        Před vykreslením stránky v prohlížeči, se~nejprve vytvoří DOM, který obsahuje všechny objekty (HTML prvky), které se~mají zobrazit.
        K~nim pak prohlížeč přiřadí CSS styly a~vykreslí je~na~obrazovku. K~DOMu pak můžeme přistupovat pomocí Javascriptu
        a~můžeme aktivně měnit obsah strány, včetně stylů nebo sledovat uživatelské akce a~na~jejich základě spouštět vlastní funkce.
        
        \subsection{HTML s dynamickým obsahem}
        HTML stránky mohou mít i~dynamický obsah. To znamená, že~se~mohou měnit podle toho, jak s~nimi uživatel pracuje.
        Na úrovni serveru se~může generovat HTML stránka, která obsahuje dynamický obsah. V~dnešní době se~častěji používá Javascript,
        nicméně v minulosti se~používaly i~jiné technologie, jako např.~PHP, které na~serveru vygenerovaly HTML stránku.

        To samozřejmě znamenalo, že~při~každé změně obsahu bylo potřeba stránku znovu načíst. Tento problém právě adresuje Javascript.
        Odesílá požadavky na~server a~na~základě odpovědi mění obsah celé nebo části stránky.

        To~znamená, že~když např.~uživatel klikne na~tlačítko, nemusí se~stránka znovu načítat celá, pouze se~načtou data ze~serveru
        a~přidají se~do~stránky. To~je~mnohem rychlejší a~příjemnější pro~uživatele.
        
        \subsection{Client-Server}
        Vývoj webových aplikací je~založen na~modelu klient-server. Klient může být webový prohlížeč, ale i~jiná, např. mobilní aplikace.
        \textbf{Server} je~počítač, který poskytuje služby klientovi. V~případě webové aplikace je~to~server, který poskytuje HTML stránku, styly a~skripty.
        Taky posléze poskytuje data, která se~do~ní načítají a~ovlivňují její obsah.

        \textbf{Klient} je~počítač, který služby serveru využívá. V~případě webové aplikace je~to~webový prohlížeč, který zobrazuje HTML stránku.
        V~případě mobilní aplikace je to~aplikace, která se~spouští na~mobilním zařízení a~komunikuje se~serverem. Např.~messengery, sociální sítě.

        \subsection{Webová služba, webový server}
        Webová služba je~program, který běží na~serveru a~poskytuje data klientovi. Na serveru může běžet více webových služeb,
        rozlišují se portem, na~kterém běží. Např.~webová služba HTTP, která poskytuje webová data a~soubory bude běžet na~portu~80,
        resp.~443 pro~HTTPS. Další služby mohou být např.~emailová služba SMTP, která poběží nejčastěji na~portu~25 a další.
        Služby \uv{naslouchají} na svých portech a~přijímají požadavky od~klientů.

        \subsection{Webový server}
        Webový server se~stará o~odpovědi většinou na požadavky, většinou přicházející přes \emph{HTTPS}, dříve \emph{HTTP}.
        Server dostane požadavek ve~formě URL\footnote{Uniform Resource Locator, někdy označované URI = Uniform Resource Identifier} adresy,
        která obsahuje cestu k~souboru, který se~má vrátit.

        Webový server dnes ale~neslouží pouze k~poskytování statických souborů, ale i~k~poskytování dynamických dat.
        Může například zajistit, že~se~před odesláním soubor na~serveru zpracuje (např. pomocí PHP) a~do~něj se~vloží aktuální data.
        Nebo může před odesláním provést kompresi dat. To~znamená, že~data se~před odesláním zmenší, čímž se~odesílají rychleji a~nezatěžují
        tak~síť. Taky to uspoří mobilní data a~klient, který přistupuje přes mobilní síť, vám bude vděčný. Na~straně klienta se~pak data rozbalí a~zobrazí.

        Webový server je~také zodpovědný za~směrování doménových jmen ke~konkrétním datům na~daném serveru. To~znamená, že~na~jednom serveru může běžet více webových stránek.
        Pomocí metody~\emph{Proxy~pass} můžeme také přesměrovat celý provoz domény na jiný port (např.~na~naši aplikaci).

        Příkladem webového serveru může být Apache, Nginx, Lighttpd, Caddy ale taky IIS\footnote{Internetová Informační Služba/Internet Information Service}\dots

        
        \subsection{MVC - Model View Controller}
        Nejvýznamnějším bodem vývoje bylo zavedení architektury MVC. MVC je~zkratka pro~\textbf{Model View Controller}.
        Ten odděluje logiku aplikace od~jejího zobrazení. To~umožňuje vývoj funkční části aplikace nezávisle na~jejím zobrazení.

        Díky tomu můžeme mít jednu aplikaci, která bude mít různé zobrazení, ale čerpá pouze z~jednoho zdroje. Např.~webovou stránku, mobilní aplikaci a desktopovou aplikaci.

        Má tři složky; \textbf{Model}, který se~stará o~manipulaci s~daty, \textbf{Controller}, který říká \textbf{Modelu}, co má dělat
        a~\textbf{View}, který se~stará o~výsledné zobrazení dat.

        Taky to umožňuje vývoj aplikace týmem vývojářů, kde každý tým může pracovat na~jiné části aplikace.

        \subsection{Frontend, Backend}
        Aplikace se~díky architektuře \emph{MVC} běžně dělí na~\textbf{frontend} a~\textbf{backend}, které pracují spolu, ale jsou oddělené.
        \textbf{Frontend} je~část aplikace, kterou vidí uživatel. Může to být např.~prohlížeč. Stará se~o interakci s uživatelem a~komunikuje s~\textbf{backendem},
        zejména pomocí asynchronních požadavků na~server. Je~napsán většinou v~HTML, CSS a~Javascriptu.

        \textbf{Backend} je~část aplikace, která běží na~serveru a~zpracovává nebo~poskytuje data \textbf{frontendu}.
        Spolupracuje s~databází a~dalšími službami, které jsou potřeba pro~správné fungování aplikace.
        \textbf{Backend} může být napsán v~různých jazycích, např.~PHP, Java, C\#, Go, Javascript, C++\dots


        \subsection{Verzovací systém GIT}
        Verzovací systém je~program, který slouží ke~správě verzí zdrojových souborů.
        Umožňuje ukládat historii změn a~vracet se~k~předchozím verzím. Taky mohou pracovat s~jednotlivými větvemi,
        které mohou obsahovat různé změny, které lze následně sloučit do~jedné větve.
        To~umožňuje vývojářům pracovat na~různých problémech projektu paralelně a~následně je~sloučit do~jednoho celku.

        \textbf{GIT} vytvořil Linus Torvalds v~roce 2005. Je~to~distribuovaný verzovací systém. To znamená, že~každý vývojář má~lokální kopii celého projektu,
        tedy může pracovat i~bez~připojení k~internetu a~následně své~změny nahrát na~server, když se~připojí.

        \subsection{Linux a příkazová řádka, WSL2}
        Linux je~operační systém, který vytvořil Linus Torvalds v~roce 1991. Je~to~nejrozšířenější open source operační systém,
        používaný zejména na~serverech. Je~to~také základ pro~mnoho dalších operačních systémů, jako např.~Android.
        
        Velikou výhodou je otevřenost a~možnost upravovat zdrojový kód. To~umožňuje vývojářům upravovat systém podle svých potřeb.
        Další výhodou je~jeho bezpečnost. V~porovnání s~ostatními operačními systémy je~mnohem bezpečnější a~méně náchylný k~virům.

        Nabízí mnoho grafických prostředí, takže si~uživatel může vybrat, jakým způsobem bude systém používat a je možné téměř
        všechno nastavit tak, jak v práci potřebuje.

        Navíc většina nástrojů pro vývoj webových aplikací je~vyvíjena pro~Linux, takže je~možné je~používat bez~problémů a~automaticky
        vše nastavit skriptem z~příkazové řádky.

        Na Windows je možné Linux využívat pomocí WSL2 (Windows Subsystem for Linux). Jedná se~o~virtuální stroj, který běží na~pozadí
        a umožní uživateli Windows používat výhody, které nabízí Linux.

        \subsection{JSON - Javascript Object Notation}
        JSON je~formát pro~výměnu dat. Je~to~textový formát, to znamená, že je~čitelný pro~člověka i~stroj.
        Používá se~ke~strukturovanému ukládání a~přenosu dat. Dříve byl~hojně využíván formát XML,
        který je~podobný struktuře HTML. Ostatně XML stále používá spousta aplikací, zejména státní správy.

        Formát JSON je~podobný zápisu objektů v~Javascriptu. Obsahuje dvojice klíč-hodnota, které jsou odděleny dvojtečkou.
        Klíčem je~řetězec a~hodnota může být číselná, řetězcová, logická hodnota, pole nebo objekt.
        Práce s~formátem JSON je~pohodlnější než s~formátem XML, protože je~mnohem jednodušší a~čitelnější a vyžaduje méně kódu,
        jak na straně serveru, tak na straně klienta a~jde~celkově o~úspornější formát.

        \subsection{API - Application Interface}
        API je~rozhraní, které umožňuje komunikaci mezi aplikacemi. V~případě webových aplikací se~jedná o~rozhraní mezi frontendem a~backendem.
        Některá rozhrání umožňují taky komunikaci na~stejném zařízení nebo mezi aplikacemi na~různých zařízeních.
        Někerá API jsou veřejná a~umožňují komunikaci s~aplikacemi třetích stran. Např.~API Google Maps.

        V~dnešní době jsou nejrozšířenější formou API \emph{REST} a~\emph{GraphQL}.

        Pro komunikaci mezi interními aplikacemi se~často používá \emph{SOAP}. Různý hardware používá k~ovládání např.~\emph{UART} a~\emph{SPI}.
        Velice oblíbený je~také výkonný přístup gRPC\footnote{Google Remote Procedure Call}, který je~založen na~protokolu HTTP/2 a~používá binární formát \emph{Protocol Buffers}.
        RPC je~zkratka pro~\emph{Remote Procedure Call} a znamená, že pouze volá vzdálené funkce, které jsou na~serveru,
        což je mnohdy efektivnější přístup.
        
        Další výhodou, že~gRPC zastává deklarativní přístup, kde definujeme služby a~jejich rozhraní a~nástroj posléze vygeneruje kód pro~jazyk, který používáme.

        \subsection{REST - Representational state transfer}
        \textbf{REST} je~architektura, která definuje způsob komunikace mezi klientem a~serverem.
        V dnešní době je~díky své jednoduchosti a~rozšířenosti nejčastěji používanou architekturou pro~tvorbu webových služeb.
        Jednoduchost spočívá zejména v~tom, že~každý \emph{Endpoint}\footnote{Endpoint/Koncový bod je~adresa zdroje, kam posíláme požadavek}
        má svou vlastní adresu.
        \textbf{REST} je~bezstavový, takže každý požadavek je~nezávislý na~předchozím, server si~neukládá data v~závislosti na~připojeném klientovi.
        To~umožňuje jednoduché testování a~ladění aplikace. Také to~ale~přináší některé výzvy a~omezení, které je~potřeba řešit.
        \textbf{REST} definuje, jaké HTTP metody se~používají pro~komunikaci, a~každé přiděluje určitý význam.
    
        Mezi nejčastěji používané metody patří GET, POST, PUT, DELETE a PATCH. Metoda je~uložena v~hlavičce HTTP požadavku.

        Popis funkce přenosu dat po síti pomocí HTTP metod je~následující:
        \begin{itemize}
            \item \textbf{GET} slouží pro~získání dat ze~serveru,
            \item \textbf{POST} pro~odeslání dat na~server,
            \item \textbf{PUT} pro~úpravu dat na~serveru,
            \item \textbf{DELETE} slouží pro~smazání dat ze~serveru,
            \item \textbf{PATCH} potom pro~úpravu části dat na~serveru.
        \end{itemize}

        Souhrnně těmto operacím říkáme \emph{CRUD operace}, což je~zkratka pro~\textbf{Create}, \textbf{Read}, \textbf{Update} a \textbf{Delete}.

        Kromě požadavků \textbf{GET} se v~těle požadavku posílají data v~textové formě, např.~formátu~JSON. 
        V~případě požadavku \textbf{GET} se~data vkládají do URL adresy, nejsou šifrována a~mají limitovanou velikost maximální délkou \emph{URL adresy}.

        \subsection{AJAX - Asynchronous Javascript Request}
        AJAX je~technologie, která umožňuje asynchronní komunikaci mezi frontendem a~backendem.
        To~znamená, že~při~každém požadavku na~server se~stránka nemusí znovu načítat, ale pouze se~načtou data ze~serveru a~přidají se~do~stránky.

        Funkce Javascriptu, umožňující asynchronní požadavky na~server, jsou \texttt{fetch} a~\texttt{XMLHttpRequest}.

        AJAX využívá různé HTTP metody, jako GET, POST, PUT, DELETE, PATCH. Tyto metody se~používají pro~odlišení typu požadavku.

        \subsection{Zabezpečení}
        Vystavením aplikace na~internetu uvádíme naše API do~nebezpečí útoku a~prolomení zabezpečení. Proto je~potřeba aplikaci zabezpečit.
        To~znamená, že~je~potřeba zabezpečit API, databázi a~další služby, které aplikace využívá. Je~nutné taky zvážit,
        jaké služby potřebujeme veřejně dostupné a~zbytek zablokovat pomocí firewallu, případně ke~službám, ke~kterým potřebujeme soukromý přístup,
        přistupovat pomocí tzv.~SSH tunelu.

        Z~hlediska zabezpečení webové aplikace pak~hovoříme o~dvou základních bodech. \textbf{Autentikace} a~\textbf{autorizace}.

            \subsubsection{Autentikace (Authentication)}
            \textbf{Autentikace} je~proces ověření identity uživatele. To~znamená, že~uživatel musí dokázat, že~je~tím, za~koho se~vydává.
            To~se~dělá pomocí hesla, které uživatel zná, nebo pomocí certifikátu, který uživatel má.

            \textbf{Autentikací} nazýváme také proces, kdy ověřujeme, že požadavek na~server přišel z~naší aplikace a~ne z~jiného zdroje,
            např.~od~útočníka. Na~to~existují různé metody, např.~využití \emph{JWT}.

            Tzn.~do~aplikace vložíme tajný klíč, který aplikace použije k~autentikaci a po~úspěšné autentikaci aplikace se~vygeneruje textový řetězec\footnote{tzv.~Token},
            který se~při~každém požadavku na~server posílá v~hlavičce požadavku.

            Při autentikaci aplikace je důležité zkontrolovat také, že~požadavek přišel z~povolené webové adresy, geografickou polohu klientovy IP adresy
            nebo další informace, které by~mohly odhalit nepovolený přístup.

            \subsubsection{Autorizace (Authorization)}
            Autorizace je~proces ověření oprávnění uživatele. To~znamená, že~uživatel musí dokázat, že~má~povolení k~provedení akce.
            To~je možné ověřit pomocí rolí, které uživatel má, nebo pomocí oprávnění k~operacím s~daným zdrojem.

            Třeba, zda je~uživatel autorem příspěvku a~tedy jej může upravit, pozměnit nebo smazat.

            \subsubsection{JWT - JSON Web Token}
            Metoda zabezpečení, která se~často používá, je~JWT. Jedná se~o~formát textového řetězce (tokenu), který obsahuje informace o~uživateli.
            Při autentikaci se~vygeneruje token\footnote{řetězec znaků}, který se~při~každém požadavku na~server posílá v~hlavičce požadavku.
            Data v~tokenu jsou zašifrována a~podepsána privátním šifrovacím klíčem, takže je~nelze změnit. V~případě změny dat v~tokenu se~token stává neplatným.
            To~umožňuje ověřit, že~token byl vytvořen serverem a~ne někým jiným. Zároveň v~datech tokenu můžeme mít informace o~uživateli, který je právě přihlášen,
            a~další kontrolní informace, jako např.~čas platnosti tokenu nebo IP adresa, ze~které byl požadavek na vytvoření tokenu prvně odeslán.

            Výhodou je také možnost přenositelnosti. Token může být vytvořen na~jednom serveru a~ověřen na~jiném, který ná k~dispozici daný klíč.

        \subsection{Údržba aplikace}
        Údržba aplikace je~důležitá část vývoje. Je~potřeba aplikaci pravidelně aktualizovat a~opravovat chyby, které se~vyskytnou.
        Tedy jde o nejdelší část životního cyklu aplikace. V~této části je~potřeba aplikaci testovat a~opravovat chyby, které se~vyskytnou.
        
            \subsubsection{SDLC - Životní cyklus vývoje software (Software Development Life Cycle)}
            SDLC je~životní cyklus vývoje softwaru. Jedná se~o~soubor procesů, které se~používají při~vývoji softwaru.
            Tyto procesy se~opakují v~kolech, které se~nazývají iterace. V~každé iteraci se~vyvíjí část aplikace, která je~následně testována.
            Vývojáři se~také mohou v~každé iteraci vrátit k~předchozí části a~upravit ji~podle potřeby.

            Existuje několik modelů \texttt{SDLC}, které se~liší počtem iterací a~způsobem vývoje. Nejčastěji se~používá model \textbf{Agile},
            který je~založen na~iteracích a~pravidelné komunikaci se~zákazníkem. V~každé iteraci se~vyvíjí část aplikace, která je~následně testována.

            Existují ale i~jiné modely vývoje, jako např.~\textbf{Vodopádový}, který je~založen na~jedné iteraci.
            V~té se~vyvíjí celá aplikace a~až poté se~testuje a~nasazuje.
            V~případě chyby je~potřeba se~vrátit na~začátek a~vyvíjet aplikaci znovu.

            \subsubsection{Sledování chyb nahlášených uživateli}
            Pokud uživatel narazí na~chyby v~aplikaci, měl~by~je~mít~možnost nahlásit. Vývojář by~potom měl mít možnost tyto chyby sledovat a~opravit~je.

            Podobný je~i~případ požadavků na~nové funkce. Uživatel by~měl mít~možnost požádat o~novou funkci nebo zlepšení stávajících,
            aby~vývojář mohl aplikaci dále rozvíjet.

            \subsubsection{Běhové a kompilační chyby}
            Kompilace je~proces překladu zdrojového kódu do~binárního souboru, který je~přímo spustitelný na~procesoru.
            Při kompilaci se~kontroluje správnost kódu, datových typů ap. Případné chyby překladač oznámí vývojáři.
            Některé chyby může překladač opravit automaticky a~podle závažnosti chyby a~námi zavedených parametrů pro~překlad
            překlad proběhne v~pořádku nebo se~zastaví a~je~potřeba chyby nejdříve opravit.

            Chyby se~ale mohou projevit až~při~běhu programu. To~jsou chyby, které např.~závisí na~datech, která se~do~programu načítají.
            Příkladem takové chyby může být špatná práce s~pamětí, která může vést k~přetečení paměti a~následnému pádu programu.

            Na běhových chybách se~docela často zakládají útoky na~aplikace. Útočník může záměrně zasílat data, která způsobí chybu v~aplikaci,
            čímž může získat přístup k~systému nebo získat data, ke~kterým by~jinak přístup neměl.

            Důležité pro opravu chyb je~uvědomit si~jejich příčinu. To~je~často složité, protože chyba se~může projevit až~po~několika krocích nebo pouze za~specifického stavu aplikace.

            Proto je~nutné aplikaci testovat a~opravovat chyby, které se~vyskytnou. To~je~důležitá část vývoje aplikace. Důležité je~také sledovat
            chyby, které se~vyskytnou uživatelům. Jeden z~nástrojů, které to~umožňují, je~Sentry.

            \subsubsection{Sentry}
            Sentry je~nástroj, který umožňuje sledovat chyby za~běhu aplikace.
            Pomocí webového rozhrání potom může vývojář (nebo jeho tým) sledovat chyby, které se~vyskytly za~chodu~aplikace.
            Sentry umožňuje sledovat chyby v~různých jazycích, jako v~našem případě \emph{Javascript a Go}. Ale i~v~dalších.


            V~každém z~těchto jazyků umožní \emph{Sentry} sledovat chyby specifické pro~daný jazyk a~tím zjednodušit opravu chyb.
            Např.~v~případě \emph{Javascriptu} umožní sledovat akce, které uživatel provedl, před tím, než došlo k~chybě. 
            Toho využiji pro~sledování chyb v~backendu i~frontendu aplikace.

            Díky tomu má~vývojář k~dispozici spoustu informací, které potřebuje k~nalezení, vyvolání a~opravě chyby.
            Sentry také umožňuje přímo vytvářet \emph{úkoly}, na kterých mohou vývojáři pracovat.
            To~umožňuje vývojářům efektivněji rozdělovat práci na~opravách chyb, které Sentry zachytí.

            Informace navíc zůstanou uloženy v~Sentry a vývojáři se~k~nim mohou kdykoliv vrátit.
            Sentry je možno integrovat do~různých nástrojů, jako např.~Slack, který se~běžně používá ke~komunikaci v~týmech
            a~upomínka na~chybu se~objeví přímo v~komunikačním kanálu týmu.
	
	\section{Návrh aplikace}
        \subsection{UCD - Use Case Diagram}
        \subsection{ERD - Entity Relationship Diagram}
        \subsection{Návrh struktury databáze}
        \subsection{Výběr technologií}
            Výběr technologií je~důležitý krok. Je nutné vybírat efektivní technologie, které budou dostatečně výkonné a~zároveň budou
            přívětivé pro vývojáře, aby vývoj nebyl zbytečně časově náročný.

            \textbf{React} je~knihovna pro~tvorbu uživatelského rozhraní. Je~vyvíjena společností Facebook a~je~používána při~vývoji mnoha webových aplikacích.
            Je oblíbená pro svou jednoduchost a~rychlost. Je~také dobře dokumentovaná a~má~velkou komunitu, která ji~podporuje.
            Existuje mnoho nástrojů, které usnadňují vývoj aplikací v~Reactu, jako
            např.~\emph{Create React App}\footnote{Někdy se setkáme se zkratkou CRA} nebo \emph{Vite}.
            Také je spousta již hotových komponent, které můžeme v~naší aplikaci používat.

            Jednotlivé komponenty mohou mít svůj vlastní stav, který se~může měnit. To~umožňuje vytvářet interaktivní aplikace.
            Komponenty mohou být funkcionalní nebo třídní. Funkcionální komponenty jsou jednodušší a~rychlejší na~vytvoření.
            Třídní komponenty jsou složitější, ale umožňují využít všech výhod Reactu.

            React je~díky využití Virtuálního DOMu velice rychlý. Virtuální DOM je~kopie DOMu, která se~používá pro~porovnání změn.
            To~umožňuje aplikaci rychleji reagovat na~změny a~měnit pouze ty části, které změnily svůj stav.

            Samozřejmě, že je možné React používat i~bez~těchto nástrojů, ale tyto nástroje usnadňují vývoj a~jsou dnes již standardem.

            \textbf{Vite} je~nástroj pro~rychlé nastartování projektu v~Reactu ale~i~v~Javascriptu, Typescriptu, Vue a dalších frameworků.
            Vite je~vyvíjen společností \emph{Vue.js}, která je~známá díky svému frameworku \emph{Vue.js}, který je~Reactu podobný.
            Také zrychluje spouštění aplikace a~nastavuje hot reload, což znamená, že~při~každé změně kódu se~aplikace v~prohlížeči automaticky načte znovu.
            Vite se~také chlubí, že~dokáže spouštět aplikace rychleji než jiné nástroje.

            \textbf{Typescript} samozřejmě použijeme i~při vývoji aplikace v~Reactu. \emph{Facebook} s~tím~samozřejmě počítá a~vydává
            kromě Reactu také balíček s~příslušnými datovými typy, pro použití v~Typescriptu.
            
            \textbf{MathJax} je~knihovna, která umožňuje zobrazovat matematické vzorce na~webových stránkách.
            Přijímá formáty matematiky zapsané v~\LaTeX{u} nebo MathML a~převádí je~na~formu, kterou může zobrazit prohlížeč.
            Díky tomu můžeme v~aplikaci názorně zobrazit matematické vzorce, které se~vyskytují v~příkladech nebo~učebních materiálech,
            jako bychom je psali v~\LaTeX{u}.

            \textbf{Go} je staticky typovaný, kompilovaný jazyk vyvinutý společností Google v~roce 2007. Je to~jazyk, který je~velice rychlý a~jednoduchý na~psaní.
            Protože je kompilovaný, je~rychlý i~při~běhu. Je~to~jazyk, který je~velice oblíbený pro~tvorbu webových aplikací.
            Umožňuje při~vývoji jednoduše importovat existující balíčky z~GIT repozitářů nebo i~jiných zdrojů a~používat je~v~aplikaci a~sám
            si~tyto balíčky stáhne a~nainstaluje. Tím~značně urychluje vývoj aplikace.

            Zároveň při kompilaci vzniká jeden binární soubor, takže přestože je aplikace docela rozsáhlá, je~dostatečně rychlá.
            
            V aplikaci \textbf{Go} používám pro~tvorbu backendu, který bude zpracovávat požadavky od~klienta a~komunikovat s~databází.

            Při~vývoji také využiju balík \textbf{cosmtrek/Air} z~Githubu, který umožňuje automaticky restartovat aplikaci při~každé změně kódu.
            Tak bude možné vývoj aplikace trochu zjednodušit.

            \textbf{Gin} je~webový framework pro~jazyk Go. Je~jednoduchý a~rychlý. Umožňuje vytvářet jednotlivé endpointy,
            které budou zpracovávat požadavky od~klienta.

            Každý \emph{koncový bod} může mít svou vlastní \emph{Obslužnou rutinu} (\textbf{handler}), která zpracuje požadavek a~vrátí odpověď.
            Gin také umožňuje vytvářet skupiny \emph{koncových bodů}, které mají společný základ~URL adresy, které mohou mít společnou \emph{obslužnou rutinu}.
            To~umožňuje vytvářet aplikace, které mají jednotnou strukturu a~jsou přehledné. Např.~všechny \emph{koncové body} pro~přihlášené
            uživatele můžeme předem definovat a~zabezpečit. 
            
            \textbf{GORM} je~ORM\footnote{ORM - Object-relational mapper} pro~jazyk Go. Umožňuje vytvářet databázové modely ze~struktur v~jazyce Go, 
            ze kterých se~automaticky vytvoří všechny databázové tabulky. Změny v~databázových modelech se~automaticky promítnou do~databáze při~migraci modelů.
            Přidání nových řádků stejně jako jejich úprava je~potom jednoduché.

            Pokud bychom potřebovali nějakou složitější operaci, samozřejmě můžeme využít i~samostatně formulované SQL dotazy.
            
            \textbf{MySQL} je~relační databáze, která je~dnes jednou z~nejpoužívanějších databází. Je~dostatečně rychlá a~spolehlivá.

            \textbf{Docker, docker-compose} jsou nástroje, které umožňují vytvářet kontejnery, ve~kterých můžeme spouštět aplikace.
            Použití kontejnerů při vývoji je~výhodné, protože vždy, když si~projekt znovu stáhneme, máme jistotu, že~aplikace běží
            ve~stejném prostředí, jako kdekoliv jinde. To~umožňuje jednodušší vývoj a~nasazení aplikace.

            Taky nám to~umožní rychle vytvořit kontejnery s~potřebnými službami, jako třeba databázi. Navíc je~kontejner odstíněn od síťového provozu,
            takže je~o~něco bezpečnější než klasická instalace databáze na~serveru.

            Na serveru navíc lze provozovat více kontejnerů, dostupných na~různých portech, což umožňuje provozovat více aplikací na~jednom serveru.

            \textbf{Nginx} se~definuje jako webový server s~možností reverzního proxy serveru a~load balanceru.
            \emph{Reverzní server} umožňuje přesměrovat požadavky na~různé servery nebo aplikace (porty), podle zadaných pravidel.
            Také je velmi rychlý ve srovnání s~jinými servery, jako např.~Apache. Zejména pak při obsluze statických souborů.

            A~nakonec \textbf{Alpine Linux}, který je~odlehčenou distribucí Linuxu, která je~určena pro~nasazení v~kontejnerech,
            nicméně je~možné ji~použít i~na~serveru. Celková velikost distribuce je~pouze 5~MB, což je~velice málo.

            Tuto distribuci používám jako hlavní systém na~svém VPS\footnote{VPS - Virtual Private Server}.
            Na~serveru poběží pouze \textbf{Nginx}, který nasměruje požadavky na~správné porty,
            \textbf{Docker}, který bude spouštět kontejnery s~aplikací a~databází a \emph{Firewall} \textbf{UFW}, který zablokuje přístup ke všem portům,
            kromě těch, které jsou potřeba pro chod aplikace.

            Tím bych chtěl dosáhnout maximálního poměru výkonu a~bezpečnosti.
	
	\section{Postup implementace}
        \subsection{Backend}
        \subsection{Frontend}
        \subsection{Testování}
        \subsection{Údržba}
        
	\section{Budoucí rozvoj aplikace}
        \subsection{Vývoj}
        Rád bych do~budoucna rozšířil aplikaci o~další funkce, jako např.~možnost vytvářet vlastní učební materiály.

        \subsection{Integrace do jiných prostředí}
        \subsection{Režim cvičení}
        \subsection{Výukové materiály a minikvízy}
        \subsection{Online zkoušení a výuka}
	
	\section{Závěr}

\printbibliography
\end{document}
